\section{Machine Learning for Cross-Sectional Returns}

\citet{gu2020empirical} provide the canonical empirical benchmark for
machine-learning-based cross-sectional asset pricing.
Using a large equity-characteristic panel, they show that tree ensembles
and neural networks outperform linear baselines when feature sets are rich
and model capacity is regularized appropriately.
\citet{freyberger2020dissecting} reinforce this conclusion from a
nonparametric perspective, demonstrating that the incremental predictive
power of asset characteristics often lies in nonlinear interactions rather
than in any single characteristic taken in isolation.
Viewed jointly, this literature suggests that the central modeling problem
in cross-sectional return prediction is not merely temporal extrapolation,
but the extraction of structured relative information from a heterogeneous
feature set.

The cryptocurrency literature confirms that the asset class is predictable,
but it is less settled on how much model complexity is actually useful.
\citet{liu2022cryptocurrency} establish the basic risk-factor structure of
cryptocurrency returns, documenting momentum and size premia that resemble
equity analogues only superficially.
More directly related to our setting, \citet{cakici2022crypto_ml} show that
machine-learning models can generate economically meaningful gains in the
cross-section of cryptocurrency returns, although the incremental benefit of
very high model complexity is limited.
\citet{filippou2024crypto_ml} reach a complementary conclusion: with a rich
set of network, momentum, technical, and attention variables, nonlinear
methods deliver substantial out-of-sample predictive and economic value.
These studies jointly suggest that the key question is not whether
cryptocurrency returns are predictable, but how the learning problem should
be formulated so that heterogeneous predictors can be turned into a stable
cross-sectional ranking.

\section{Learning to Rank in Finance}

The learning-to-rank literature \citep{cao2007learning} offers a natural
foundation for long-short portfolio formation because the investment problem
is inherently ordinal: the objective is to place assets in the correct
relative order, not necessarily to forecast their exact subsequent return.
\citet{poh2021building} provide the closest precursor to our approach.
They show that standard ListNet training becomes unstable when raw returns
are used directly as softmax targets, because target sharpness varies across
weeks with the cross-sectional dispersion of returns.
Replacing raw-return targets with cross-sectional rank percentiles produces
substantially more stable long-short performance in equities.

Our contribution builds directly on this insight, but in a setting where the
problem is arguably more severe.
Cross-sectional volatility heterogeneity is considerably larger in
cryptocurrency markets than in equities, so any objective that depends on
raw-return scale is more prone to unstable gradients.
We therefore adopt the rank-percentile ListNet target and complement it with
gradient accumulation, treating stable objective design as a first-order
methodological requirement rather than a minor training detail.

\section{Temporal Deep Learning Architectures}

Long Short-Term Memory networks \citep{hochreiter1997long,fischer2018lstm_finance}
remain a strong baseline for financial sequence modeling because they can
extract local temporal dependence with relatively controlled parameter count.
This is especially relevant when predictors exhibit short-lived persistence,
as in event-driven social media activity or ETF flow sequences.

The Temporal Fusion Transformer \citep{lim2021temporal} extends this design
space by combining recurrent encoding, Variable Selection Networks (VSN),
and attention-based temporal aggregation.
Relative to generic Transformer architectures \citep{vaswani2017attention},
TFT is particularly attractive for financial panels with heterogeneous
predictors because it provides an explicit mechanism for feature selection
and interpretability.
At the same time, its larger capacity also makes it more vulnerable to
training instability when the effective sample size is limited.
This trade-off is one motivation for our side-by-side comparison of LSTM,
TFT, and a purely cross-sectional architecture.

\section{Alternative Data Sources}

The use of alternative data in financial prediction is now well established.
\citet{tetlock2007giving} show that media tone predicts short-horizon stock
returns, while \citet{bollen2011twitter} document predictive content in
Twitter-derived mood measures.
In the cryptocurrency setting, \citet{maitre2025social} provide more direct
evidence that social-media-based attention is related to the cross-section of
crypto returns.
These studies collectively motivate the use of text-derived or attention-like
signals as complements to price and fundamental information.

Our Trump social media variables differ from general sentiment proxies in an
important respect: they are designed to capture policy-linked signals from a
single high-impact political actor rather than diffuse market mood.
This makes them structurally closer to event variables than to broad
attention measures.

ETF flow data provides a second form of alternative information that is
distinct from both prices and on-chain activity.
Since the approval of spot Bitcoin ETFs in January 2024, the daily flow data
published by Farside Investors
\citep{bitcoin_etf_farside,ethereum_etf_farside} has made it possible to separate
legacy-holder rotations out of Grayscale products from genuinely new capital
entering the asset class.
That decomposition is economically meaningful because it distinguishes
supply-side pressure from demand-side expansion, and is therefore a
potentially useful predictor of cross-sectional dispersion even if it is not
always informative about market direction.

\section{Statistical Inference}

Evaluating financial machine-learning strategies requires inference tools
that remain valid under serial dependence and non-normal return
distributions.
We therefore use Newey--West heteroskedasticity and autocorrelation
consistent statistics \citep{newey1987simple,newey1994automatic} and the
stationary bootstrap of \citet{politis1994stationary}.
These methods do not improve predictive performance, but they reduce the
risk of mistaking concentrated short-sample performance for robust evidence,
which is particularly important given the 49-week test window in this study.
